% !TEX root = ../bb_AiRa_Kappaleet.tex

% ö = \%c3\%b6
% ä = \%C3\%A4

\xdef\sectiontitle{\IIIsectiontitle} %aiheuttama
\TOCrefs

%%%%%%%%%%%%%%%
\begin{frame}[label=3_4_a]%[label=3_4_TOC1]
\frametitle{\texorpdfstring{\culine[\sectiontitleulinecolor]{\sectiontitle}}{\sectiontitle} }
\label{\TOCname}

\vspace{-0.5cm}\label{\TOCname} % vertical spacing is off, 0.7cm doesn't work

\begin{itemize}[leftmargin=0.5cm,itemsep=\chapterTOCitemsep]
    \item[3.1]<1-> \hyperlink{3_1_a}{\IIIaSubsectiontitle}
    \item[3.2]<1-> \hyperlink{3_2_a}{\IIIbSubsectiontitle}	
    \item[3.3]<1-> \hyperlink{3_3_a}{\IIIcSubsectiontitle}	
    \item[3.4]<1-> \hyperlink{3_4_a}{\CDAlert[red]{\IIIdSubsectiontitle}}
    \item[3.5]<1-> \hyperlink{3_5_a}{\IIIeSubsectiontitle}
    \item[3.6]<1-> \hyperlink{3_6_a}{\IIIfSubsectiontitle}
    \item[3.7]<1-> \hyperlink{3_7_a}{\IIIgSubsectiontitle}
\end{itemize}

%\endofslide 
\end{frame}
%%%%%%%%%%%%%%%

\xdef\sectiontitle{\IIIdSubsectiontitle} 

%%%%%%%%%%%%%%%
\begin{frame}
\frametitle{\texorpdfstring{\culine[\sectiontitleulinecolor]{\sectiontitle}}{\sectiontitle} }
\framesubtitle{Radioaktiiviset hajoamisprosessit}%Radioactive Decay processes
\appear{}

\semismall
%On Radioactive Decay
\begin{itemize}
	\item Noin \CDAlert[fuchsia]{3000 nuklidia on löydetty}\appear{, \CDAlert[red]{suurin osa} niistä \CDAlert[red]{on radioaktiivisia}} 
	
	\item \CDAlert[fuchsia]{Epästabiileilla} radioaktiivisilla nuklideilla \appear{on taipumus spontaanisti emittoida pieniä energeettisiä hiukkasia} 
	\item Arvioiden mukaan\appear{, jo tunnettujen lisäksi} \appear{noin 2000 on vielä löytämättä}\appear{, jotka kaikki ovat epävakaita ja radioaktiivisia}
	
	\item  On olemassa \CDAlert[blue]{useita 'hajoamisprosesseja'}, esim.\
\begin{itemize}
	\item $\alpha$-hajoaminen ($\mathrm{He}^{2+}$-ioni)
	\item $\beta^{-}$-emissio (elektroni)\appear{, ja  $\beta^{+}$-emissio} \appear{(positroni)}
	\item $\gamma$-säteily \appear{(energeettistä $1 \mathrm{\;keV}–10 \mathrm{\;MeV}$ sähkömagneettista säteilyä)}
\end{itemize}

\item[] Ja  \CDAlert[tuniblue]{vähemmän tunnettuja} prosesseja:
\begin{itemize}
\item Protoniemissio
\item Neutroniemissio
\item Positroniemissio
\item Kaksoisbeeta-emissio (kaksi elektronia)
\item Näiden kombinaatioita 
\end{itemize}

\end{itemize}

\endofslide 
%\endofsection
\end{frame}
%%%%%%%%%%%%%%%

%%%%%%%%%%%%%%%
\begin{frame}
\frametitle{\texorpdfstring{\culine[\sectiontitleulinecolor]{\sectiontitle}}{\sectiontitle} }
\framesubtitle{Radioaktiiviset hajoamisprosessit}%Radioactive Decay processes
\appear{}
% 6.5 cm = midway, 10 cm = 3/4 
%\vspace{2.5mm}
\begin{minipage}{8.5cm}
\begin{itemize}
	\item \Alert[red]{Semi-empiiristä sidosenergian kaavaa} \appear{voidaan käyttää arvioimaan isotooppien stabiilisuutta}
	\appear{\xdef\loc{south east}
\begin{tikzpicture}[remember picture,overlay] % anchor=south
    \node (A) [yshift=-0.35*\figYshift,xshift=\figXshift, anchor=\loc] at (current page.\loc) {\includegraphics[height=8.5cm]{Figures_booklet/SOM_Tipler_Nuclear_physics_figures/SOM_Tipler_Nuclear_physics_figure_20.png}}; % 8cm max height
\end{tikzpicture}\CR[ref = t, pos=left]{}}
	\item Kokeiden varmistamana, me voimme hahmotella sidosenergioita\appear{ (ks. kuva missä} \appear{z-akselia on skaalattu tekijällä  $100 \mathrm{\;MeV}$)} 
	
	\item Kaavan avulla voidaan myös arvioida \CDAlert[blue]{raja-arvoja} sille kuinka monta protonia\appear{/neutronia ytimessä voi olla} 
	\implies{ Nämä arviot $N-Z$ rajoille \\
		muodostavat  alemman kuvan siniset ääriviivat}
\end{itemize}
\end{minipage}



\endofslide 
%\endofsection
\end{frame}
%%%%%%%%%%%%%%%


%%%%%%%%%%%%%%%
\begin{frame}
\frametitle{\texorpdfstring{\culine[\sectiontitleulinecolor]{\sectiontitle}}{\sectiontitle} }
\framesubtitle{Radioaktiiviset hajoamisprosessit}%Radioactive Decay processes

\seminormal
% 6.5 cm = midway, 10 cm = 3/4 
%\vspace{2.5mm}
\begin{minipage}{8.5cm}
\begin{itemize}
	\item \appear{\xdef\loc{south east}%
	\begin{tikzpicture}[remember picture,overlay] % anchor=south
    \node (A) [yshift=-0.35*\figYshift,xshift=\figXshift, anchor=\loc] at (current page.\loc) {\includegraphics[height=8.5cm]{Figures_booklet/SOM_Tipler_Nuclear_physics_figures/SOM_Tipler_Nuclear_physics_figure_20.png}};% 8cm max height
\end{tikzpicture}\CR[ref = t, pos=left]{}}%
Radioaktiiviset prosessit noudattavat fysikaalisia \CDAlert[magenta]{säilymislakeja}, joissa säilyvät suureet ovat:
\begin{itemize}
	\item Relativistinen massa–energia
	\item Sähköinen varaus 
	\item Liikemäärä
	\item Kulmaliikemäärä
\end{itemize}

\item Ja myös:
\begin{itemize}
	\item Nukleoniluku
	\item Leptoniluku
\end{itemize}

\item Jotta radioaktiivista isotooppia löytyisi luonnosta \appear{(Maasta)}\appear{, sen \Alert[red]{puoliintumisajan} täytyy olla pitempi kuin Maan iän} \appear{($4,5\times 10^{9}$ vuotta)}\appear{, tai nuklidia täytyy \CDAlert[blue]{jatkuvasti syntyä}} \appear{radioaktiivisen nuklidin hajoamisen seurauksena}
\end{itemize}
\end{minipage}

\endofslide 
%\endofsection
\end{frame}
%%%%%%%%%%%%%%%

%%%%%%%%%%%%%%%
\begin{frame}
\frametitle{\texorpdfstring{\culine[\sectiontitleulinecolor]{\sectiontitle}}{\sectiontitle} }
\framesubtitle{Alfahajoaminen}%Alpha decay
%
\appear{}
 
\semismall
% 6.5 cm = midway, 10 cm = 3/4 
%\vspace{2.5mm}
\begin{minipage}{9.0cm}
	\begin{itemize}
	\item Alfahajoaminen on \Alert[red]{energeettisen alfa-hiukkasen} \appear{(eli He-4 ytimen)} \appear{spontaania emissiota}
\appear{\xdef\loc{north east}%
\begin{tikzpicture}[remember picture,overlay]% anchor=south
    \node (A) [yshift=1*\figYshift,xshift=\figXshift, anchor=\loc] at (current page.\loc) {\includegraphics[width=4.7cm]{Figures_booklet/SOM_12_Nuclear_physics_figures/SOM_Nuclear_physics_Figure_18.png}}; % 8cm max height
\end{tikzpicture}\CR[ref = h and t]{}}%
\appear{Hajoamisen aikana alfahiukkanen emittoituu ytimestä} \appear{tunneloitumalla potentiaalivallin läpi}
\appear{\xdef\loc{south east}%
\begin{tikzpicture}[remember picture,overlay] % anchor=south
    \node (A) [yshift=-0.4*\figYshift,xshift=\figXshift, anchor=\loc] at (current page.\loc) {\includegraphics[width=4.7cm]{Figures_booklet/SOM_Tipler_Nuclear_physics_figures/SOM_Tipler_Nuclear_physics_figure_1_cropped.png}};% 8cm max height
\end{tikzpicture}}

\item Radioaktiiviselle hajoamiselle on aina syy.
\appear{Että ydin voisi ylipäätänsä olla radioaktiivinen}\appear{, sen \Alert[blue]{massan täytyy olla suurempi kuin sen hajoamistuotteiden massojen}.} \appear{Vapautunut energia on tällöin}
\[
Q  \equal \appear{\textrm{((} m_{i}}  \minus \appear{m_{f}-}\appear{m_{\alpha} \textrm{))}} \appear{ c^{2}}
\]

\item $\alpha$-aktiivinen \CDAlert[fuchsia]{emonuklidi} on painava \appear{(yleensä $A>200)$}\appear{, jolloin myös \CDAlert[fuchsia]{tytärnuklidi} on painava}\appear{, kun taas alfahiukkanen on kevyt: }\appear{\quad (\CDAlert[fuchsia]{liikemäärän säilyminen})}
\[
\appear{m \vec{v_t}}  \plus  \appear{M \vec{V_e}} \equal \appear{0} \quad \Rightarrow \quad \appear{v_t=}\frac{M}{m} \appear{V_e} \quad  \Rightarrow \quad  \frac{1}{2} \appear{m v_t^{2}}  \equal \frac{M}{m} \frac{1}{2} \appear{M V_e^{2}}
\]

\end{itemize}
\end{minipage}
%\vspace{2.5mm}

\endofslide 
%\endofsection
\end{frame}
%%%%%%%%%%%%%%%

%%%%%%%%%%%%%%%
\begin{frame}
\frametitle{\texorpdfstring{\culine[\sectiontitleulinecolor]{\sectiontitle}}{\sectiontitle} }
\framesubtitle{Alfahajoaminen}%Alpha decay
\appear{\xdef\loc{north east}
\begin{tikzpicture}[remember picture,overlay] % anchor=south
    \node (A) [yshift=1*\figYshift,xshift=\figXshift, anchor=\loc] at (current page.\loc) {\includegraphics[width=4.7cm]{Figures_booklet/SOM_Tipler_Nuclear_physics_figures/SOM_Tipler_Nuclear_physics_figure_25.png}}; % 8cm max height
\end{tikzpicture}\CR[ref = h and t]{}}%
% 6.5 cm = midway, 10 cm = 3/4 
%\vspace{2.5mm}
\begin{minipage}{8.8cm}
\begin{itemize}
	\item Nuklideja sitoo toisiinsa vahva voima. \appear{Kaukana ytimestä voima kuitenkin heikkenee}\appear{, koska vahvan vuorovaikutuksen kantama on lyhyt ($\sim$fm).} \appear{Suurilla etäisyyksillä potentiaalivalli muodostuu siis \CDAlert[blue]{positiivisesta Coulombisesta potentiaalista}}
	
	\item Oletuksia alfahajoamisen ymmärtämiseksi:
\end{itemize}
\end{minipage}
%\vspace{2.5mm}

\begin{itemize}
\item<.->[] \begin{itemize} \small
		\item Alfahiukkanen voi ennestään olla olemassa painavassa ytimessä 
	\item Ytimen potentiaalikaivossa \appear{poukkoileva alfahiukkanen on jatkuvassa liikkeessä} 
	\item On olemassa pieni todennäköisyys sille  \appear{että hiukkanen tunneloituu vallin läpi} \appear{sen törmätessä valliin}
\end{itemize}

\item Itse asiassa on myös ehdotettu\appear{ (ainakin kevyiden ydinten tapauksissa)}\appear{, että nuklidit koostuisivat alfahiukkasista...}
\end{itemize}

\endofslide 
%\endofsection
\end{frame}
%%%%%%%%%%%%%%%

%%%%%%%%%%%%%%%
\begin{frame}
\frametitle{\texorpdfstring{\culine[\sectiontitleulinecolor]{\sectiontitle}}{\sectiontitle} }
\framesubtitle{Alfahajoaminen ja hajoamisketjut}%Alpha decay
\appear{}

\begin{itemize}
	\item \CDAlert[red]{Kaikki painavat ytimet $(A > 210)$} \appear{\CDAlert[red]{ovat epävakaita} koska emoytimen massa on suurempi} \appear{kuin sen hajoamistuotteiden} \appear{(tytärnuklidin ja alfahiukkasen)} \appear{yhteenlaskettu massa} 

\item $\alpha$-emissiossa ydin muuttuu: \appear{$\Delta Z=\Delta N=-2$} \appear{eli $\Delta A=-4$} 
\item Prosessit voivat myös ketjuuntua muodostaen \CDAlert[blue]{hajoamissarjoja} \vspace{2mm}
\[ 
\footnotesize
\begin{array}{r | c | c | c | c}
& \text {A} & \text {Pitkäikäisin isotooppi} & t_{1/2} & \text {Vakaa tytär} \\
\hline Toriumsarja & 4 n  	   & n=58  & 1,4 \times 10^{10} \;a & Pb-208 \\
 Neptuniumsarja & 4 n+1 & n=58 & 2,1 \times 10^{6} \;a & Tl-205\\
 \text{Uraani- eli Radiumsarja} & 4 n+2 & n=59 & 4,5 \times 10^{9} \;a & Pb-206\\
 Aktiniumsarja & 4 n+3 & n=58 & 7,0 \times 10^{8} \;a & Pb-207
\end{array}
\vspace{2mm}
\]

\item Hajoamissarjoissa \appear{tapahtuu \CDAlert[blue]{peräkkäisiä radioaktiivisia hajoamisia}.} \appear{Jotta hajoamissarjat voisivat seurailla stabiilisuusviivaa}\appear{, \CDAlert[red]{myös beetahajoamista täytyy tapahtua}} \appear{(kasvattamaan järjestyslukua $Z$} \appear{tai parantamaan $N/Z$-suhdelukua)}

\end{itemize}

\endofslide 
%\endofsection
\end{frame}
%%%%%%%%%%%%%%%

%%%%%%%%%%%%%%%
\begin{frame}
\frametitle{\texorpdfstring{\culine[\sectiontitleulinecolor]{\sectiontitle}}{\sectiontitle} }
\framesubtitle{Stabiilisuusviiva ja hajoamisketjut}%Line of stability and decay chains

\begin{minipage}{9.0cm}
\begin{itemize}
	\item Yhtä lukuunottamatta sarjat \CDAlert[fuchsia]{löytyvät myös} \CDAlert[fuchsia]{luonnosta}.\appear{ Luonnossa neptunium\-sar\-jan (Np-237)}\appear{ neptunium-atomit ovat jo ehtineet hajota muiksi nuklideiksi}
\appear{\xdef\loc{north east}
\begin{tikzpicture}[remember picture,overlay] % anchor=south
    \node (A) [yshift=1*\figYshift,xshift=\figXshift, anchor=\loc] at (current page.\loc) {\includegraphics[width=4.7cm]{Figures_booklet/SOM_Tipler_Nuclear_physics_figures/SOM_Tipler_Nuclear_physics_figure_26.png}}; % 8cm max height
\end{tikzpicture}\CR[ref= t]{}}%
\item Kuvassa on toriumsarjan hajoamisketju %\appear{and Uranium series}

\item \CDAlert[blue]{Stabiilisuusviiva} on merkitty katkoviivalla. \appear{Tytärnuklidien päätyessä kauas viivan vasemmalle puolelle} \appear{, eli \CDAlert[fuchsia]{neutronirikkaalle puolelle}}\appear{, prosessi etenee beetahajoamisten kautta tasapainottaen $N/Z$-suhdelukua}

\end{itemize}
\end{minipage}
\vspace{2.5mm}

\begin{itemize}
\item Huoneilmassa esiintyvä \CDAlert[red]{radonkaasu} \appear{(tyypillistä Suomessa)} \appear{on Rn-222}\appear{, mikä kuuluu radium}\appear{/uraanisarjaan (U-238).} \appear{Sen puoliintumisaika on neljä päivää}\appear{,  mutta sitä muodostuu hajoamissarjan seurauksena jatkuvasti}
\end{itemize}

\endofslide 
%\endofsection
\end{frame}
%%%%%%%%%%%%%%%

%%%%%%%%%%%%%%%
\begin{frame}
\frametitle{\texorpdfstring{\culine[\sectiontitleulinecolor]{\sectiontitle}}{\sectiontitle} }
\framesubtitle{Hajoamisketjut}
\appear{}
\seminormal
% 6.5 cm = midway, 10 cm = 3/4 
%\vspace{2.5mm}
\begin{minipage}{9.0cm}
\begin{itemize}
	\item \CDAlert[blue]{Toriumsarja}: \appear{Kuvassa näkyy lukuisia Radium-223:n energiatasoja} \appear{joihin Th-227 voi alfahajoamisen seurauksena muuntua}
	\appear{\xdef\loc{north east}%
\begin{tikzpicture}[remember picture,overlay] % anchor=south
    \node (A) [yshift=1*\figYshift,xshift=\figXshift, anchor=\loc] at (current page.\loc) {\includegraphics[width=4.7cm]{Figures_booklet/SOM_Tipler_Nuclear_physics_figures/SOM_Tipler_Nuclear_physics_figure_29.png}}; % 8cm max height
\end{tikzpicture}\footnote{\TiplerCRshort}}
\item Th-227:sta emittoituneen $\alpha$-säteilyn spektri on alla olevassa kuvassa \appear{(Ra-223:n energia (keV) alaindeksinä)}
\appear{\xdef\loc{south west}%
\begin{tikzpicture}[remember picture,overlay] % anchor=south
    \node (A) [yshift=-0.1*\figYshift,xshift=-\figXshift, anchor=\loc] at (current page.\loc) {\includegraphics[width=8cm]{Figures_booklet/SOM_Tipler_Nuclear_physics_figures/SOM_Tipler_Nuclear_physics_figure_27.png}}; % 8cm max height
\end{tikzpicture}}
\item Huom! Usein tytärnuklidi Radium-223 päätyy virittyneeseen tilaan. Toisinaan\appear{ ylimääräinen energia emittoituu lisäksi $\gamma$-säteilynä}
\end{itemize}
\end{minipage}
%\vspace{2.5mm}



\endofslide 
%\endofsection
\end{frame}
%%%%%%%%%%%%%%%

%%%%%%%%%%%%%%%
\begin{frame}
\frametitle{\texorpdfstring{\culine[\sectiontitleulinecolor]{\sectiontitle}}{\sectiontitle} }
\framesubtitle{Beetahajoaminen}%Beta decay
\appear{}

% 
\begin{itemize}
	\item Sähköstaattinen vetovoima ei voi vangita elektronia ytimeen\appear{, mutta tästä huolimatta elektroneitakin emittoituu välillä ytimistä.}\appear{ Nähtävästi hajoamisprosesseissa syntyy elektroneita}

\item Huom! \CDAlert[red]{$\beta$-säteilyssä}\appear{ \CDAlert[red]{neutroni muuttuu protoniksi}}\appear{, jolloin tytärnuklidi on eri alkuaine kuin emonuklidi.} \appear{Protonin ja elektronin} 
\end{itemize}

% 6.5 cm = midway, 10 cm = 3/4 
%\vspace{2.5mm}
\begin{minipage}{8.2cm}
\begin{itemize}
	\item[] emittoituessa\appear{, varaus säilyy.} \appear{Energioille:}
$$
\begin{aligned}
\appear{Q} & \equal  \appear{\textrm{((} m_{n}-m_{p}-m_{e} \textrm{))}}  \appear{c^{2}} \\
& \equal \appear{(939,57-938,28-0,511) \mathrm{\;MeV}}  \equal \appear{0,78 \mathrm{\;MeV}}
\end{aligned}
$$
\appear{Oletetaan tämän olevan syntyneiden hiukkasten kineettistä energiaa}\appear{, enimmäkseen elektronin}\appear{, koska $m_{p} / m_{e}=1840$.} \appear{Kokeellisesti on elektronien energioiden on havaittu vaihtelevan välillä $0-Q$}
\end{itemize}
\end{minipage}
%\vspace{2.5mm}

\appear{\xdef\loc{south east}
\begin{tikzpicture}[remember picture,overlay] % anchor=south
    \node (A) [yshift=-0.25*\figYshift,xshift=\figXshift, anchor=\loc] at (current page.\loc) {\includegraphics[width=5.3cm]{Figures_booklet/SOM_12_Nuclear_physics_figures/SOM_Nuclear_physics_Figure_19.png}}; % 8cm max height
\end{tikzpicture}
\footnote{\HarrisCRshort}}

\endofslide 
%\endofsection
\end{frame}
%%%%%%%%%%%%%%%

%%%%%%%%%%%%%%%
\begin{frame}
\frametitle{\texorpdfstring{\culine[\sectiontitleulinecolor]{\sectiontitle}}{\sectiontitle} }
\framesubtitle{Beetahajoaminen}%Beta decay
\appear{}

%\seminormal
\begin{itemize}
	\item Historiallisesti oli olemassa kolme ongelmaa:
\begin{itemize}
	\item Havainto elektronien kineettisen energian vaihtelusta välillä 0  \appear{(ja \Alert[red]{lähes $0,78 \mathrm{\;MeV}$})}\appear{, oli mahdotonta selittää vain kahden lopputuotteen avulla} \appear{(tytär ja elektroni)}
	\appear{\xdef\loc{south east}
\begin{tikzpicture}[remember picture,overlay] % anchor=south
    \node (A) [yshift=-0.25*\figYshift,xshift=\figXshift, anchor=\loc] at (current page.\loc) {\includegraphics[width=5.3cm]{Figures_booklet/SOM_12_Nuclear_physics_figures/SOM_Nuclear_physics_Figure_20.png}}; % 8cm max height
\end{tikzpicture}
\footnote{\HarrisCRshort}}
\end{itemize}
\end{itemize}

% 6.5 cm = midway, 10 cm = 3/4 
\vspace{2.8mm}
\begin{minipage}{8.5cm}
\begin{itemize}
\item<.->[] \begin{itemize}[itemsep=2.8mm]
	\item n\appear{, p} \appear{ja e$^-$}\appear{, ovat 1/2-spinin fermioneita:}
	$$ 
	\begin{array}{cccccc}
	 & \appear{n} & \rightarrow & \appear{p} 	   &  \plus  & \appear{e^-} \\
	 \appear{\text {\small \CDAlert[blue]{spin} }} & \appear{\text {\small \Alert[blue]{1/2}}} & & \appear{\text {\small \Alert[blue]{1/2}} }&  &  \appear{\text {\small \Alert[blue]{1/2}}}
	\end{array} \qquad \quad
	$$
	\appear{eli lopputuotteiden kokonaisspinin täytyy olla kokonaisluku.} \appear{\CDAlert[blue]{Eikö spin säily} \Alert[blue]{prosessissa?}}
	
\item e$^-$:n energia \appear{on myös toisinaan lähes $Q$}\appear{, joten kolmannen lopputuotteen}\appear{ \CDAlert[fuchsia]{massan} täytyisi olla \CDAlert[fuchsia]{erittäin pieni}}
\end{itemize}
\end{itemize}
\end{minipage}
%\vspace{2.5mm}



\endofslide 
%\endofsection
\end{frame}
%%%%%%%%%%%%%%%

%%%%%%%%%%%%%%%
\begin{frame}
\frametitle{\texorpdfstring{\culine[\sectiontitleulinecolor]{\sectiontitle}}{\sectiontitle} }
\framesubtitle{Beetahajoaminen ja neutriino}%Beta decay and neutrino
\appear{}
%\seminormal
\begin{itemize}
	\item Wolfgang Pauli \appear{ja Enrico Fermi} \appear{ehdottivat että olisi olemassa} \appear{kolmas lopputuote}\appear{, \CDAlert[red]{neutriino}}\appear{, jonka:}
	\appear{\xdef\loc{south east}%
\begin{tikzpicture}[remember picture,overlay] % anchor=south
    \node (A) [yshift=-0.25*\figYshift,xshift=\figXshift, anchor=\loc] at (current page.\loc) {\includegraphics[width=5.3cm]{Figures_booklet/SOM_12_Nuclear_physics_figures/SOM_Nuclear_physics_Figure_20.png}};% 8cm max height
\end{tikzpicture}\footnote{\HarrisCRshort}}%\CR{}
\begin{itemize}
	\item spin olisi $1 / 2$

\item massa olisi lähes nolla \appear{\textrm{((}ehkä n.\ $\sim 2,2 \mathrm{\;eV} / \mathrm{c}^{2} \textrm{))} $}

\item vuorovaikutus muun aineen kanssa olisi \\
	erittäin heikkoa  \implies{ niinsanottu \CDAlert[fuchsia]{heikko vuorovaikutus} }
\end{itemize}
\item Itse asiassa muodostunutta hiukkasta $\bar{\nu}_{e}$ \\ \appear{kutsutaan  nykyään \CDAlert[blue]{antineutriinoksi}}
\end{itemize}

\endofslide 
%\endofsection
\end{frame}
%%%%%%%%%%%%%%%

%%%%%%%%%%%%%%%
\begin{frame}
\frametitle{\texorpdfstring{\culine[\sectiontitleulinecolor]{\sectiontitle}}{\sectiontitle} }
\framesubtitle{$\beta^{-}$-hajoamisen esimerkki (boori-12)}
\appear{}

%
\begin{itemize}
	\item Tarkastellaan boori-12:ta beetahajoamista:
\[
{ }_{\;5}^{12} B \quad \rightarrow \quad \appear{{ }_{\;6}^{12} C}  ~\plus~ \appear{{ }_{\,-1}^{~0} \beta^{-}}  ~\plus~ \appear{\bar{\nu}_{e}}
\]

\item Neutroni muuttuu protoniksi\appear{, syntynyt kevyempi ydin} \appear{(tytärnuklidi)} \appear{on tiukemmin sidottu kun $Z=N$} \appear{(ja molemmat ovat kokonaislukuja).} \appear{Eli kun boori hajoaa, se madaltaa energiaansa}

\item Lasketaan energiat \appear{(unohdetaan $\beta^{-}$:n massa}\appear{, koska hiilen ${ }_{\;6}^{12} C$ atomimassa} \appear{sisältää jo yhden ylimääräisen elektronin):}
$$
\begin{aligned}
\appear{Q} & \equal  \appear{\textrm{((} m_{\text {emo}}} \minus \appear{m_{\text {tytär}}}  \minus \appear{m_{\beta} \textrm{))}}  \appear{c^{2}}  \equal  \appear{\textrm{((} m_{B}-m_{C} \textrm{))}}  \appear{c^{2}} \\
&  \equal \appear{ \textrm{((} 12,014352 \;u-12 \;u \textrm{))}} \appear{\times 931,5 \;MeV / u} \equal \appear{13,4 \mathrm{\;MeV}}
\end{aligned}
$$
\end{itemize}

\endofslide 
%\endofsection
\end{frame}
%%%%%%%%%%%%%%%

%%%%%%%%%%%%%%%
\begin{frame}
\frametitle{\texorpdfstring{\culine[\sectiontitleulinecolor]{\sectiontitle}}{\sectiontitle} }
\framesubtitle{Typpi-12 beetahajoaminen}
\appear{}
\seminormal
\begin{itemize}
	\item Tarkastellaan nyt typpi-12:ta \appear{(${ }_{\,7}^{12} N$)} \appear{$\beta^{+}$-hajoamista}\appear{, jolloin protoni muuttuu neutroniksi:}%
	\appear{\xdef\loc{south east}%
\begin{tikzpicture}[remember picture,overlay] % anchor=south
    \node (A) [yshift=-0.35*\figYshift,xshift=\figXshift, anchor=\loc] at (current page.\loc) {\includegraphics[width=4.6cm]{Figures_booklet/SOM_12_Nuclear_physics_figures/SOM_Nuclear_physics_Figure_21.png}}; % 8cm max height
\end{tikzpicture}\CR{}}%
\[
{ }_{\,7}^{12} N \quad \rightarrow \quad \appear{{ }_{\;6}^{12} C}  ~\plus~ \appear{{ }_{+1}^{\,0} \beta^{+}} \appear{~+~ \nu}
\]
\item<.->[] \[ %what's this here
p^{+} \quad \rightarrow \quad \appear{n}  ~\plus~  \appear{e^{+}} \appear{~+~\nu}
\]
\appear{missä $\beta^{+}$ on positiivinen elektroni}\appear{, positroni}\appear{, mikä varmistaa sekä varauksen että kulmaliikemäärän (spinin) säilymisen }
\end{itemize}

% 6.5 cm = midway, 10 cm = 3/4 
\vspace{2.5mm}
\begin{minipage}{8.5cm}
\begin{itemize}
\item Lasketaan taas reaktion $Q$:
$$
\begin{gathered}
\appear{Q = \textrm{((} m_{\text {emo}}-m_{\text {tytär}}-m_{\beta} \textrm{))}  c^{2}= \textrm{((} m_{N}-m_{C}-2 m_{e} \textrm{))}  c^{2}} \\
\appear{= \textrm{((} 12,018613 \;u-12 \;u-2 \cdot 5,486 \times 10^{-5} \;u \textrm{))}  c^{2} =16,3 \mathrm{\;MeV}}
\end{gathered}
$$

\item Huom: \Alert[blue]{vapaa protoni ei voi muuttua neut\-roniksi ja positroniksi}, \appear{mutta ytimeen sidottu protoni voi}
\end{itemize}

\end{minipage}
%\vspace{2.5mm}

\endofslide 
%\endofsection
\end{frame}
%%%%%%%%%%%%%%%

%%%%%%%%%%%%%%%
\begin{frame}
\frametitle{\texorpdfstring{\culine[\sectiontitleulinecolor]{\sectiontitle}}{\sectiontitle} }
\framesubtitle{Elektronikaappaus}
\appear{}
%Electron capture
\begin{itemize}
	\item\appear{\xdef\loc{south east}%
\begin{tikzpicture}[remember picture,overlay] % anchor=south
    \node (A) [yshift=-0.32*\figYshift,xshift=\figXshift, anchor=\loc] at (current page.\loc) {\includegraphics[width=4.6cm]{Figures_booklet/SOM_12_Nuclear_physics_figures/SOM_Nuclear_physics_Figure_21.png}};% 8cm max height
\end{tikzpicture}\CR{}}%
\CDAlert[red]{Elektronikaappausreaktiossa} \appear{ytimen protoni kaappaa elektronin joltain atomin elektronikuorelta} \appear{ja muuttuu neutroniksi} 
	
	\item \CDAlert[red]{Järjestysluku $Z$ pienenee yhdellä ja $N$ kasvaa}, kuten $\beta^{+}$-hajoamisessa. \appear{Monet $\beta^{+}$-aktiiviset ytimet}\appear{ voivat kaapata elektroneita}

\item Aina kun järjestysluvulla $Z$ olevan ytimen massa on pienempi luvulla $Z-1$ olevan\appear{, voi tapahtua elektronikaappaus}
\end{itemize}

% 6.5 cm = midway, 10 cm = 3/4 
\vspace{2.5mm}
\begin{minipage}{8.5cm}
\begin{itemize}
	\item Jos massa-ero on vähintään $2 m_{e} c^{2}$\appear{, niin ydin voi olla myös $\beta^{+}$-aktiivinen.} \appear{Tällöin nämä kaksi \CDAlert[blue]{hajoamisprosessia kilpailevat}}

\item Elektronikaappauksen todennäköisyys on\\ 
	usein erittäin pieni. \appear{Paitsi jos atomin elektroni on ytimen välittömässä läheisyydessä}
\end{itemize}
\end{minipage}
%\vspace{2.5mm}

\endofslide 
%\endofsection
\end{frame}
%%%%%%%%%%%%%%%

%%%%%%%%%%%%%%%
\begin{frame}
\frametitle{\texorpdfstring{\culine[\sectiontitleulinecolor]{\sectiontitle}}{\sectiontitle} }
\framesubtitle{Esimerkki elektronikaappauksesta}
\appear{}

% 
\begin{itemize}
	\item Tarkastellaan nyt hiili-11:sta \appear{(${ }_{\,6}^{11} C$)} \appear{tapahtuvaa elektronikaappausta }
\[
{ }_{\,6}^{11} C \appear{~+~} \appear{ _{-1}^{\;0} \beta^{-}} \quad \rightarrow \quad \appear{{ }_{\,5}^{11} B}  ~\plus~ \appear{\nu} \qquad \appear{\textrm{((}} \appear{p^{+}}  ~\plus~ \appear{e^{-}} \quad \rightarrow \quad  \appear{n~+~\nu} \appear{\textrm{))}}
\]
\appear{Vapautunut energia} \appear{riippuu ytimen ja tyttären massaerosta (eli \CDAlert[red]{massavajeesta}):} 
\[
Q  \equal  \appear{\textrm{((} m_{\text {emo}}}  \minus \appear{m_{\text {tytär}} \textrm{))} } \appear{c^{2}}
\]

\item Toinen tyypillinen esimerkki elektronikaappauksesta on 
\[
 { }_{24}^{51} C r  ~\plus~ \appear{ _{-1}^{\;0} \beta^{-}} \quad \rightarrow \quad \appear{{ }_{23}^{51} V}  ~\plus~ \appear{\nu_{e}} \appear{\,,~ \text{missä}} \qquad  \appear{Q=}\appear{0,751 \mathrm{\;MeV}}
 \]
\end{itemize}

\endofslide 
%\endofsection
\end{frame}
%%%%%%%%%%%%%%%

%%%%%%%%%%%%%%%
\begin{frame}
\frametitle{\texorpdfstring{\culine[\sectiontitleulinecolor]{\sectiontitle}}{\sectiontitle} }
\framesubtitle{Hajoamisketjut}%Decay chains
\appear{}

\semismall
% 6.5 cm = midway, 10 cm = 3/4 
%\vspace{2.5mm}
\begin{minipage}{9.0cm}
\begin{itemize}[itemsep=1.5mm]
	\item Tarkastellaan vielä lisää torium- ja uraanisarjoissa tapahtuvia alfahajoamisia $(Z,N)$-kuvaajan avulla
\appear{\xdef\loc{east}
\begin{tikzpicture}[remember picture,overlay] % anchor=south
    \node (A) [yshift=0.*\figYshift,xshift=\figXshift, anchor=\loc] at (current page.\loc) {\includegraphics[height=8cm]{Figures_booklet/SOM_12_Nuclear_physics_figures/SOM_Nuclear_physics_Figure_23.png}}; % 8cm max height
\end{tikzpicture}
\footnote{\HarrisCRshort}}
\item Jotta isotoopit pysyisivät \CDAlert[fuchsia]{stabiilisuusviivan} lähellä\appear{, täytyy tapahtua myös beetahajoamisia (kun $N \gtrsim Z$).}

\item Lopulta \CDAlert[red]{sarjat päättyvät} lähelle \Alert[red]{lyijyn} isotooppeja: \appear{Pb-208 (torium)}\appear{, Pb-206 (uraani)} \appear{ja Pb-207 (aktinium).} \appear{Nämä isotoopit eivät ole erityisen vakaita}\appear{, vaan sattuvat olemaan  \CDAlert[red]{'ensimmäiset'} vakaat isotoopit}\appear{, päättäen hajoamissarjat.} \appear{(toisin kuin eräät kevyet isotoopit kuten K-40, C-14 jne.)}

\item Huomaa myös:
\begin{itemize}
	\item $\beta^{-}$-hajoaminen tapahtuu \appear{diagonaalisesti \\alas} \appear{\& oikealle }
	\item $\beta^{+}$-hajoaminen \appear{tapahtuu diagonaalisesti \\ylös} \appear{\& vasemmalle}
\end{itemize}

\end{itemize}
\end{minipage}
%\vspace{2.5mm}

\endofslide 
%\endofsection
\end{frame}
%%%%%%%%%%%%%%%

%%%%%%%%%%%%%%%
\begin{frame}
\frametitle{\texorpdfstring{\culine[\sectiontitleulinecolor]{\sectiontitle}}{\sectiontitle} }
\framesubtitle{Stabiilisuuden laakso}%Valley of stability
\appear{}

\semismall
% 6.5 cm = midway, 10 cm = 3/4 
%\vspace{2.5mm}
%\begin{minipage}{8.0cm}
\begin{itemize}[itemsep=1.6mm]
	\item Beetahajoamisten avulla \appear{tutkitaan myös ydinten} \appear{\CDAlert[blue]{stabiilisuuden laaksoa}}

\item Beetahajoamisten avulla \appear{voidaan liikkua kohtisuorasti energialaakson pitkään akseliin nähden}

\item Samanpainoisten ydinten energiaprofiilit näyttävät hahmottelevan parabolisen energialaakson (kuva a)%
\xdef\loc{south east}%
\begin{tikzpicture}[remember picture,overlay] % anchor=south
    \node (A) [yshift=-0.15*\figYshift,xshift=\figXshift, anchor=\loc] at (current page.\loc) {\includegraphics[width=11cm]{Figures_booklet/SOM_Tipler_Nuclear_physics_figures/SOM_Tipler_Nuclear_physics_figure_33.png}}; % 8cm max height
\end{tikzpicture}\footnoteleft{\TiplerCRshort}%
\item Kuvan energiat \Alert[red]{laskettiin semiempiirisen massakaavan avulla}

\item Kun $A$ on pariton\appear{, paritermi on nolla \textrm{((}$C_{5}=0$\textrm{))}}\appear{, kun $A$ on parillinen}\appear{, beetahajoamiset}
\end{itemize}
% 6.5 cm = midway, 10 cm = 3/4 
%\vspace{1.4mm}
\begin{minipage}{2.8cm}
\begin{itemize}
	\item[] hyppelevät kah-\\
		den paraabelin\\ välillä\appear{, missä $\pm C_{5}$.} \appear{(joko $Z$ = parillinen\\ ja $N$ = parillinen\\ tai $Z$ = pariton ja $N$ = pariton)}
\end{itemize}
\end{minipage}
%\vspace{2.5mm}

\endofslide 
%\endofsection
\end{frame}
%%%%%%%%%%%%%%%

%%%%%%%%%%%%%%%
\begin{frame}
\frametitle{\texorpdfstring{\culine[\sectiontitleulinecolor]{\sectiontitle}}{\sectiontitle} }
\framesubtitle{Gammahajoaminen}
\appear{}
%Gamma decay
\begin{itemize}
	\item Gammahajoamisessa%
	\appear{\xdef\loc{south east}%
\begin{tikzpicture}[remember picture,overlay] % anchor=south
    \node (A) [yshift=-0.25*\figYshift,xshift=\figXshift, anchor=\loc] at (current page.\loc) {\includegraphics[width=4.3cm]{Figures_booklet/SOM_12_Nuclear_physics_figures/SOM_Nuclear_physics_Figure_24.png}}; % 8cm max height
\end{tikzpicture}\CR{}}%
 \appear{\CDAlert[fuchsia]{virittynyt ydin hajoaa} alemmalle tilalle emittoiden fotonin}\appear{, jolloin isotooppi ei muutu.} \appear{Prosessi on analoginen virittyneen atomitilan relaksaatioprosessin kanssa, jossa myös emittoituu fotoni}

\item \CDAlert[red]{ATOMI}: Atomin prosesseissa energiat ovat suuruusluokkaa $\mathrm{\;eV}-\mathrm{\;MeV}$ \appear{ja emittoituneet aallonpituudet ovat $\mu \mathrm{m}-\mathrm{\;nm}$}

\item \CDAlert[blue]{YDIN}: Ytimen prosesseissa energiat ylittävät MeV\appear{, ja aallonpituudet ovat $\appear{\lambda} \equal \frac{h c}{E} \approx \appear{10^{-3} \mathrm{\;nm}}$
}
\item Huomaa että emittoituneen fotonin energian ollessa välillä keV–MeV\appear{, ei voida erottaa}\appear{, onko se peräisin \Alert[red]{atomin} relaksoitumisesta (\CDAlert[red]{röntgensäteilyä})} \appear{vai \Alert[blue]{ytimen} (\CDAlert[blue]{gammasäteilyä})}\appear{, jos\\ lisäinformaatiota ei ole saatavilla!}
\end{itemize}



\endofslide 
%\endofsection
\end{frame}
%%%%%%%%%%%%%%%

%%%%%%%%%%%%%%%
\begin{frame}
\frametitle{\texorpdfstring{\culine[\sectiontitleulinecolor]{\sectiontitle}}{\sectiontitle} }
\framesubtitle{Yhdistetty alfa/beeta- ja gammahajoaminen}%Combined Alpha/beta and gamma decay
\appear{}

\begin{itemize}
	\item Kun ydin on emittoinut $\alpha$-, tai $\beta$-hiukkasia\appear{,  \Alert[red]{ydin jää yleensä virittyneeseen tilaan}}\appear{, ja emittoi myös $\gamma$-säteilyä} 
	
	\item Näiden virittyneiden tilojen elinajat ovat usein erittäin lyhyitä\appear{, paitsi silloin kun valintasäännöt kieltävät virittyneen tilan purkautumisen} \appear{(vrt.\ fosforenssi)} 
	
	\item Käytännössä erisuuret $\gamma$-energiat usein \Alert[red]{liittyvät tiettyihin hajoamisprosesseihin}/siirtymiin ytimen energiatasoilla

\item Yleisesti ottaen isotooppien $\gamma$-säteilyn energiatasot ovat aika karakteristisia\appear{, ja niitä voidaan käyttää \CDAlert[red]{sormenjälkinä}}\appear{, eli tunnistamaan tiettyjä isotooppeja}\appear{, ja myös antamaan lisäinformaatiota esim. ydinten kuorien energiatasoista} \appear{(ks.\ aiempi kuva torium-227:n emissiospektristä)}
\end{itemize}

\endofslide 
%\endofsection
\end{frame}
%%%%%%%%%%%%%%%


%%%%%%%%%%%%%%%
\begin{frame}
\frametitle{\texorpdfstring{\culine[\sectiontitleulinecolor]{\sectiontitle}}{\sectiontitle} }
\framesubtitle{Sisäinen konversio}%Internal conversion
\appear{}

%Internal conversion
\begin{itemize}
	\item Virittynyt ydin voi relaksoitua gammasäteilyn emittoinnin lisäksi myös \CDAlert[blue]{sisäisen konversion} kautta

\item Konversiossa virittyneen ytimen energia \Alert[blue]{siirtyy} \CDAlert[blue]{atomin elektronille}\appear{, joka emittoituu atomista}\appear{ suurella nopeudella}
 
 \item Tällaista tapahtuu lähinnä \CDAlert[blue]{$s$-elektroneille $K$- ja $L$-kuorilla}\appear{, joilla on suurehko todennäköisyys sijaita lähellä atomin ydintä} 
 
 \item Energia on tällöin muotoa:
\[
hf  \equal \appear{ \textrm{((} E_{\text {korkea}}}\appear{-E_{\text {matala}} \textrm{))}} \appear{-E_{\text {elektroni}}}
\]

\item Sisäinen konversio on yhden vaiheen mutta usean hiukkasen prosessi\appear{, vaikka periaatteessa sen voidaan ajatella tapahtuvan kahdessa vaiheessa}\appear{, emittoiden ja absorboiden $\gamma$-fotonin}
\end{itemize}

\endofslide 
%\endofsection
\end{frame}
%%%%%%%%%%%%%%%

%%%%%%%%%%%%%%%
\begin{frame}
\frametitle{\texorpdfstring{\culine[\sectiontitleulinecolor]{\sectiontitle}}{\sectiontitle} }
\framesubtitle{Spontaani fissio}%Spontaneous fission
\seminormal
\appear{}

\begin{minipage}{8.6cm}
\begin{itemize}
	\item Jotkut painavat radioaktiiviset ytimet hajoavat vielä erilaisemmalla tavalla, \appear{eli \Alert[red]{hajoamalla kahteen suurempaan fragmenttiin} \CDAlert[red]{spontaanin fission} kautta}
\appear{\xdef\loc{north east}
\begin{tikzpicture}[remember picture,overlay] % anchor=south
    \node (A) [yshift=0.8*\figYshift,xshift=\figXshift, anchor=\loc] at (current page.\loc) {\includegraphics[width=5cm]{Figures_booklet/SOM_12_Nuclear_physics_figures/SOM_Nuclear_physics_Figure_25.png}}; % 8cm max height
\end{tikzpicture}
\footnote{\HarrisCRshort}}
\item Ajava voima on \CDAlert[fuchsia]{minimienergiaperiaate}

\item Kahdella ytimellä on  \Alert[red]{suurempi pinta-ala}\appear{ kuin alkuperäisellä ytimellä} \appear{(eli suurempi energia)} 
\end{itemize}
\end{minipage}
\vspace{2.5mm}

\begin{itemize}
\item \appear{Mutta erilliset ytimet voivat myös \Alert[blue]{järjestää protoninsa} kauemmas toisistaan} \appear{(verrattuna yksittäisen ytimen tilanteeseen)}
\implies{ \Alert[blue]{Coulombinen repulsiovoima heikkenee}\appear{, mikä kasvattaa sidosvoimia\\} \appear{ja kompensoi kasvaneen pintaenergian} }

\item Spontaanissa fissiossa on myös tärkeää muistaa\appear{, että  fragmentoituminen synnyttää myös \CDAlert[red]{vapaita neutroneita}}\appear{, mitkä voivat käynnistää uusia radioaktiivisia reaktioita törmäilemällä lähistön ytimiin}
\end{itemize}

%\endofslide 
\endofsection
\end{frame}
%%%%%%%%%%%%%%%


